\section{Docker und Kubernetes}

\subsection{Docker}
\textbf{Docker} ist eine \emph{Containerisierungs-Plattform}. Sie erlaubt es, Anwendungen zusammen mit allem, was sie zum Laufen brauchen, in sogenannte Container zu verpacken. Vor Containern gab es oft den Spruch "Works on my Machine", da ein Entwickler andere Abhängigkeiten installiert hatte als der andere Kollege. Das Problem ist durch Docker gelöst, da man alle Abhängigkeiten mitgibt.

\subsubsection{Image}
Ein \textbf{Image} ist der \textit{Blueprint}, also die Vorlage für einen Container. Es enthält alles, was zum Ausführen eines Containers erforderlich ist. Ein Image kann in einer \textit{Docker Compose} referenziert werden. In dieser werden auch noch Einstellungen vorgenommen wie Environment-Variablen.

Beispiel einer Docker Compose, mit der MySQL ausgeführt wird:

\begin{syntaxbox}[YAML]
services:
  mysql:
    image: mysql:8.0
    container_name: mysql_dev
    restart: unless-stopped

    command:
      - --default-authentication-plugin=mysql_native_password
      - --bind-address=0.0.0.0

    ports:
      - "3307:3306"

    environment:
      MYSQL_ROOT_PASSWORD: Levin-Dev-SQL

    volumes:
      - mysql_data:/var/lib/mysql
      - ./init.sql:/docker-entrypoint-initdb.d/init.sql

volumes:
  mysql_data:
\end{syntaxbox}

\subsubsection{Container}
Ausführbare Pakete mit Anwendungscode, Laufzeit, Bibliotheken und Abhängigkeiten. Sie sind \emph{zustandslos und unveränderlich}. Bei Code-Änderungen müssen sie neu erstellt werden.

\subsubsection{Container Registry}
Eine \textbf{Container Registry} ist der Ort, an dem Images gespeichert werden. Es gibt öffentliche Registries wie beispielsweise \textit{Docker Hub}, aber auch jeder Cloud-Anbieter bietet die Möglichkeit, eine eigene unternehmensinterne Container Registry zu hosten.

\subsection{Kubernetes}
\textbf{Kubernetes} oder auch \textit{K8s} wurde von Google entwickelt. Kubernetes ist eine Plattform, um Container zu \emph{orchestrieren}. Es verwaltet Container (nicht nur von Docker). Durch Kubernetes sind \emph{Hochverfügbarkeit}, \emph{Skalierbarkeit} und \emph{Disaster Recovery} möglich.

\begin{figure}[H]
    \centering
    \includegraphics[width=\textwidth]{KubernetsFeatures}
    \caption{Kubernetes Features}
\end{figure}

\subsubsection{Pods}
Ein \textbf{Pod} ist die \emph{kleinste Einheit} in Kubernetes. Ein Pod kann aus einem oder mehreren Containern bestehen. Alle Container in einem Pod haben geteilte Ressourcen wie Speicher und Netzwerk. Beispielsweise kann eine App in einem App-Pod zusammengefasst werden.

\subsubsection{Node}
Ein \textbf{Node} ist ein Server bzw. Rechner eines Clusters.

\subsubsection{Cluster}
Ein \textbf{Cluster} ist eine Sammlung von Nodes, die gemeinsam Container und Pods ausführen.

\subsubsection{Instanzen}
Eine einzelne laufende Kopie einer Ressource.

\subsubsection{Deployments}
Vergleichbar mit \textit{Docker Compose}. Ein \textbf{Deployment} enthält mehrere Container und ist sozusagen ein Bauplan für einen Pod. Allerdings wird auch angegeben, wie viele Pods existieren sollen.

\subsubsection{Services}
Ein \textbf{Service} ist die \emph{stabile Netzwerk-Abstraktion} vor einer Menge von Pods, die den Traffic an diese Pods weiterleitet. Er stellt eine stabile IP und einen Port bereit und kann Funktionen wie einen \textit{Load Balancer} konfiguriert haben, sodass die dahinterliegenden Pods erreichbar bleiben, auch wenn sie sich ändern.

\subsection{Kubernetes Funktion}
Ein Kubernetes Cluster benötigt eine \textit{Master Node}. Diese wird zur \textit{Control Plane}, die die Container auf den \textit{Worker Nodes} verwaltet. Wenn ein Node ausfällt, hat Kubernetes die \textit{Self-Healing}-Funktion und führt den Container auf einem anderen Node aus. Ein Cluster hat ein virtuelles Netzwerk, damit die Nodes untereinander kommunizieren können.